\graphicspath{{Images/}}

\section{Resumen}
Este anteproyecto propone el desarrollo de un \textbf{entorno virtual (VR)} para \textbf{entrenar a operarios} en el uso de un \textbf{brazo robótico} y la posterior \textbf{teleoperación} con un joystick diseñado ad hoc sobre \textit{Raspberry Pi Pico}. El sistema combina simulación en \textit{Unity}, medición de desempeño y un pipeline que facilite la transferencia desde el entrenamiento en VR al control del robot real.

\section{Motivación y justificación}
La teleoperación es pertinente cuando existen \textit{riesgos ambientales}, se requiere \textit{precisión manual}, hay \textit{alta variabilidad} en la tarea o se busca \textit{transferir expertise} humana. La VR reduce costos y riesgos de entrenamiento, habilita \textit{iteración rápida} y define métricas objetivas (tiempo, trayectorias, colisiones, suavidad del movimiento).

\section{Objetivos}
\subsection{Objetivo general}
Desarrollar un entorno VR y un esquema de control que permitan entrenar y luego teleoperar un brazo robótico mediante un joystick propio, priorizando seguridad, usabilidad y desempeño.

\subsection{Objetivos específicos}
\begin{enumerate}
  \item Implementar un \textbf{simulador} de brazo robótico en Unity con interacción XR.
  \item Diseñar el \textbf{mapeo de controles} del joystick (RPi Pico) a los grados de libertad del robot.
  \item Instrumentar \textbf{métricas de entrenamiento}: tiempo, errores, colisiones, suavidad/jerk.
  \item Diseñar y validar \textbf{modos de teleoperación}: velocidad limitada, asistencia, seguimiento de path.
  \item Incorporar \textbf{mecanismos de seguridad}: e-stop, zonas prohibidas, límites de torque/velocidad.
  \item Realizar una \textbf{evaluación} con usuarios: usabilidad (SUS), aprendizaje y transferencia a robot real/simulador de mayor fidelidad.
\end{enumerate}

\section{Alcance y exclusiones}
\subsection*{Alcance}
\begin{itemize}
  \item Entrenamiento VR de tareas representativas (p.ej., pick\&place, posicionamiento preciso).
  \item Teleoperación con joystick RPi Pico (prototipo funcional).
  \item Integración de comunicación (serial/UDP/TCP) y capa de seguridad básica.
\end{itemize}
\subsection*{Fuera de alcance (por ahora)}
\begin{itemize}
  \item Retroalimentación háptica avanzada.
  \item \textit{Motion planning} autónomo o visión por computadora compleja.
  \item Certificación normativa industrial.
\end{itemize}

\section{Requisitos}
\subsection{Funcionales}
\begin{itemize}
  \item Control en tiempo real del robot (simulado) desde el joystick.
  \item Escenarios VR con objetivos, \textit{checkpoints} y feedback.
  \item Registro de métricas y generación de reportes.
  \item Modo teleoperado con límites de seguridad (e-stop, zonas).
\end{itemize}
\subsection{No funcionales}
\begin{itemize}
  \item \textbf{Latencia} end-to-end $< 60$--$100$\,ms (objetivo inicial).
  \item \textbf{Usabilidad}: SUS $\geq 70$.
  \item \textbf{Compatibilidad}: Windows y \textit{headset} VR objetivo.
  \item \textbf{Mantenibilidad}: arquitectura modular y documentación.
\end{itemize}

\section{Estado del arte (breve)}
Se revisarán enfoques de simulación de robots en motores 3D, bibliotecas XR, estándares de seguridad en teleoperación y soluciones de \textit{joysticks} personalizados. 

\section{Metodología y tecnologías}
\subsection{Motor gráfico y XR}
Unity (render pipeline apropiado) y sistema de \textit{input}. 
\subsection{Arquitectura}
\begin{itemize}
  \item \textbf{Capa de Dispositivos}: joystick RPi Pico $\rightarrow$ PC (USB/serial/UDP).
  \item \textbf{Capa de Control}: mapeo joystick $\rightarrow$ comandos (velocidad/posición).
  \item \textbf{Capa de Simulación}: cinemática/dinámica y colisiones.
  \item \textbf{Capa de VR/UX}: escenas, HUD, feedback, tutoriales.
  \item \textbf{Teleoperación}: canal de comunicación + \textit{safety layer}.
\end{itemize}

\subsection{Hardware de interfaz: Joystick RPi Pico}
\begin{itemize}
  \item \textbf{Estado actual}: primer prototipo impreso 3D y firmware básico (PWM/ADC/USB HID o serial).
  \item \textbf{Plan}: iteraciones de ergonomía, \textit{deadband}, \textit{debounce}, vibración opcional.
\end{itemize}
\begin{figure}[H]
\centering
\includegraphics[angle=90,width=0.7\textwidth]{Imagenes/joystick_top_inside.png}
\end{figure}
\begin{figure}[H]
\centering
\includegraphics[angle=90,width=0.7\textwidth]{Imagenes/joystick_top_full.png}
\begin{figure}[H]
    \centering
    \includegraphics[width=0.7\linewidth]{Imagenes/gamepad_tester.png}
    \caption{Se utilizó la librería TinyUSB para implementar el protocolo HID}
    
\end{figure}
\end{figure}


\subsection{Simulación del robot}
Cinemática directa/inversa, límites articulares, velocidad y \textit{soft stops}. Opcional: exportar URDF y/o puente ROS2 para escalabilidad futura.

\subsection{Teleoperación y seguridad}
Modos de control (velocity/position), \textit{rate limiting}, filtros, geofencing, e-stop físico/virtual, \textit{watchdog} de conexión.

\section{Plan de evaluación}
\subsection{Métricas}
Tiempo por tarea, tasa de errores/colisiones, suavidad (jerk promedio), \#\ de reintentos, aprendizaje (\textit{pre/post}), SUS.
\subsection{Diseño experimental}
$N \approx 8$--$12$ participantes (piloto), ensayo de 2–3 tareas con y sin asistencia, análisis de normalidad y comparaciones (p.ej., Wilcoxon/ttest).



\section{Entregables}
\begin{itemize}
  \item Demo VR funcional (entrenamiento).
  \item Prototipo de joystick RPi Pico (STL/FW).
  \item Módulo de teleoperación con \textit{safety layer}.
  \item Informe con resultados y repositorio con código/documentación.
\end{itemize}

\section{Riesgos y mitigaciones}
\begin{itemize}
  \item Latencia alta $\rightarrow$ \textit{rate limiting}, compresión, simplificación de escena.
  \item Inestabilidad de control $\rightarrow$ saturaciones, \textit{low-pass}, límites articulares.
  \item Fatiga del usuario $\rightarrow$ sesiones cortas, ergonomía del joystick, tutoriales progresivos.
\end{itemize}

\section{Consideraciones éticas y de seguridad}
Privacidad de datos de uso, seguridad funcional (e-stop), ergonomía y prevención de mareo en VR (pausas, locomoción adecuada).

\section{Plan de bibliografía}
Se incluirán trabajos de teleoperación, entrenamiento en VR, cinemática de robots y guías de HMI/seguridad. 

